% Curs 10, pagina 1 -- manuscris datat 6-XII-1994.
% Transcrierea omite calculele intermediare tăiate în manuscris.

\clearpage
\section*{Curs 10 - Pagina 1}

\begin{center}
\underline{\textbf{ALGEBRA}}\\[2pt]
\textit{(CURS 10)}
\end{center}
\begin{flushright}
\textit{6-XII-94}
\end{flushright}

\bigskip
\noindent\underline{\textbf{Acțiunea unui grup pe o mulțime. T. lui Sylow 1, 2, 3.}}

\bigskip
\noindent\underline{Def.} Fie un grup $(G,\cdot)$, $A\neq\varnothing$, $A$ mulțime.

\noindent O aplicație $G\times A\xrightarrow{\ \psi\ } A$,
\quad $(a,x)\overset{\psi}{\longmapsto}a\cdot x\in A$, cu proprietățile:

\bigskip
\noindent (1)\quad $a\cdot(b\cdot x)=(a\cdot b)\cdot x$,
\quad $(\forall)\,a,b\in G$, \quad $(\forall)\,x\in A$.

\medskip
\noindent (2)\quad $e\cdot x=x$.

\bigskip
\noindent\underline{Ex.} \textbf{1)} Fie $G$ grup arbitrar, $A=G$,
\quad $G\times G\xrightarrow{\ \psi\ }G$.

\noindent $(a,x)\overset{\psi}{\longmapsto}a\cdot x=axa^{-1}$
(conjugatul lui $x$ prin $a$).

\bigskip
\noindent (1)\quad $a\cdot(b\cdot x)=a(bxb^{-1})a^{-1}=abxb^{-1}a^{-1}$.

\medskip
\begin{center}
\fbox{$\begin{aligned}
 a\cdot(b\cdot x)&=(ab)\cdot x,\\
 e\cdot x&=x.
\end{aligned}$}
\end{center}

\noindent $\psi$ -- acțiune prin conjugare.

\bigskip
\noindent\textbf{2)} $G$ grup, \quad $A=\{H\mid H\leq G\}$.

\noindent Definim $G\times A\xrightarrow{\ \psi\ }A$,
\quad $(g,H)\overset{\psi}{\longmapsto}g\cdot H=gHg^{-1}\in A$.

\bigskip
\noindent $f_a:G\longrightarrow G$, \quad $f_a(x)=axa^{-1}$,
\quad $f_a\in\operatorname{Aut}(G)$.

\bigskip
\noindent $f_a(xy)=axya^{-1}$.

\noindent $f_a(x)\cdot f_a(y)=axa^{-1}aya^{-1}=axya^{-1}$.

\bigskip
\noindent $f_a(x_1)=f_a(x_2)\ \Longrightarrow\ ax_1a^{-1}=ax_2a^{-1}
\ \Longrightarrow\ x_1=x_2$.

\noindent $\Longrightarrow f_a$ -- injectivă.

\medskip
\noindent $(\forall)\,y\in G$, $(\exists)\,x\in G$ a.î.\ $f_a(x)=y$;
\quad $x=a^{-1}ya$ $\Longrightarrow f_a$ -- surjectivă.

\bigskip
\noindent $\Longrightarrow f_a\in\operatorname{Aut}(G)$.
\quad $f_a$ se numește \underline{automorfism interior} al lui $G$.

\bigskip
\noindent $aHa^{-1}=f_a(H)\in A$.

\bigskip
\noindent $\psi$ satisface (1) și (2). $\psi$ se numește
\underline{acțiune prin conjugare} a lui $G$ pe mulțimea subgrupurilor sale.
